\documentclass[12pt]{article}

\usepackage[utf8]{inputenc}
\usepackage{amsmath, amssymb, amsthm}
\usepackage{geometry}
\geometry{a4paper, margin=2.5cm}

\title{Componente p-primare ale unui Modul Finit Generat}
\author{}
\date{}

\theoremstyle{definition}
\newtheorem{definition}{Definiție}

\theoremstyle{plain}
\newtheorem{theorem}{Teoremă}

\begin{document}

\maketitle

\begin{definition}
Fie $R$ un domeniu principal (PID) și $M$ un $R$-modul. Pentru un element prim $p \in R$,
definim \emph{componenta $p$-primară} a lui $M$ prin


\[
M_p = \{\, x \in M : \exists\, r \ge 1 \text{ cu } p^r x = 0 \,\}.
\]


\end{definition}

\begin{definition}
Partea torsională a lui $M$ este


\[
t(M) = \{\, x \in M : \exists\, 0 \ne r \in R \text{ cu } rx = 0 \,\}.
\]


Un modul se numește \emph{modul de torsiune} dacă $t(M) = M$, respectiv
\emph{modul liber} dacă este izomorf cu $R^k$ pentru un anumit $k$.
\end{definition}

\begin{theorem}[Decompoziția p-primară]
Fie $R$ un PID și $M$ un $R$-modul finit generat. Fie $P$ un sistem de reprezentanți
pentru clasele de asocieri ale elementelor prime din $R$.

Atunci:
\begin{enumerate}
    \item[(i)] $M_p \ne 0$ doar pentru un număr finit de primi $p \in P$;
    \item[(ii)] partea torsională se descompune ca sumă directă
    

\[
    t(M) = \bigoplus_{p \in P} M_p,
    \]


    unde suma este finită;
    \item[(iii)] în special, dacă $M$ este un modul de torsiune, atunci
    

\[
    M = M_{p_1} \oplus \cdots \oplus M_{p_s}
    \]


    pentru unici primi $p_1, \dots, p_s$ (până la unități).
\end{enumerate}
\end{theorem}

\end{document}
